\section{Introduction}\label{introduction}

Grinding is widely used as a precision finishing process when tight dimensional tolerances, controlled surface texture, and stable near-surface properties are required for high-value components. The arithmetic mean roughness \(R_a\) remains one of the most commonly reported descriptors of ground-surface quality \cite{Brinksmeier2006,Jawahir2011}, although it represents only part of the broader surface-integrity response. Surface roughness can affect fatigue performance \cite{Novovic2004}, tribological behaviour \cite{Brinksmeier2006}, and coating adhesion or tolerance compliance \cite{Jawahir2011}, particularly when the component operates under cyclic loading or sliding contact. Stylus profilometry provides accurate post-process roughness measurement, but it is typically offline, local, and too slow for immediate corrective action in a control loop. These constraints motivate in-process or near-real-time prediction of \(R_a\) from machine signals acquired during grinding.

Acoustic emission (AE), vibration, and process parameters provide complementary information on the grinding process. AE signals contain high-frequency elastic-wave components that are associated with localized contact phenomena \cite{Dornfeld1992}, such as grain--workpiece interaction, rubbing, grit fracture, wheel loading, and micro-fracture \cite{Teti2010}. Vibration signals describe the dynamic response of the wheel--workpiece--machine system, including spindle and wheel modes \cite{Altintas2004}, chatter-related behaviour \cite{Quintana2011}, and structural transmission through the machine tool \cite{Hassui2003}. Process parameters such as wheel speed, workpiece speed, and depth of cut provide additional regime-level context because the same sensor amplitude or spectral pattern may have different implications under different operating conditions. Thus, AE and vibration are physically complementary: AE is more sensitive to localized high-frequency contact phenomena, whereas vibration captures system-level dynamics. At the same time, AE interpretation must remain sensor-aware because measured high-frequency content is shaped by the sensor bandwidth, couplant, mounting path, and transfer function.

Before deep learning became common in manufacturing monitoring, roughness and process-state prediction were usually based on engineered signal descriptors. Time-domain features such as RMS, variance, skewness, and kurtosis have been widely applied to tool state classification \cite{Byrne1995}, often in combination with frequency-band energy \cite{Antoni2006} or wavelet-based descriptors and process parameters \cite{Teti2010}. These descriptors have been mapped to surface roughness using linear regression or shallow neural networks \cite{Benardos2003}, and to tool condition and chatter using support vector regression \cite{Ozel2005} or random forests and boosted-tree models \cite{Wuest2016}. Such approaches are compact, relatively data-efficient, and easier to audit than high-capacity neural networks. However, manually selected scalar features may miss local time--frequency structures, and many studies have been restricted to one signal representation or one model family, making it difficult to separate the effect of feature design from the effect of the regression model.

Recent manufacturing-monitoring studies have increasingly adopted convolutional neural networks and residual encoders \cite{Wang2018}, spectrogram-based CNNs and attention mechanisms \cite{Khan2018}, and transformers or graph-based models \cite{Lei2020}. These methods are attractive because time--frequency representations can be processed directly, allowing local spectral patterns to be learned without complete reliance on handcrafted feature definitions. Multimodal fusion can also combine AE, vibration, force, power, process parameters, and engineered descriptors at data, feature, or decision level \cite{Khaleghi2013}. When the training data are large, diverse, and representative of the future operating range, such models can provide accurate predictions. Nevertheless, their performance must be assessed under validation protocols that reflect the intended deployment scenario, particularly when the experimental dataset is small and condition-structured.

Grinding experiments are often expensive, machine-specific, and limited in the number of process conditions that can be tested. Repeated grinding passes collected under the same condition are not independent in the same sense as randomly sampled observations, because they share wheel state, process parameters, material context, machine dynamics, and measurement setup. Ordinary random train--test splits or random \(K\)-fold cross-validation can therefore leak condition identity from training to test data and yield overly optimistic estimates of generalization to unseen grinding regimes. Related concerns have been reported for cross-validation during model selection \cite{varma2006bias} and for structured datasets requiring grouped validation \cite{Roberts2017}. In addition, deep networks may be sensitive to initialization, architecture selection, validation fold, hyperparameter tuning, and post-selection when the number of independent experimental conditions is small. Their larger parameter counts and preprocessing requirements may also complicate low-latency monitoring, while learned features can be difficult to diagnose when a specific held-out condition fails. These limitations do not imply that deep learning is unsuitable for grinding monitoring; rather, they indicate that deep-learning models should be compared with transparent and more auditable baselines under identical grouped validation.

A complementary strategy is to encode domain structure before learning. Regularized linear models, such as ridge regression, provide signed coefficients after feature standardization and permit direct inspection of feature-level associations \cite{Hoerl1970}. Tree ensembles are not fully transparent, but random forests \cite{Breiman2001} and gradient-boosted trees \cite{Friedman2001} are more auditable than many deep architectures. Split statistics and permutation importance can be computed via XGBoost \cite{Chen2016} or LightGBM \cite{Ke2017}, while localized attributions are available through TreeSHAP analysis \cite{Lundberg2020}. Physics-aware representations can also reduce the burden placed on the supervised model by compressing sensor data into structured descriptors, including wavelet-scattering transforms \cite{mallat2012group} and heuristic process or sensor indicators \cite{Anden2014}. Such representations may be data-efficient and stable in small manufacturing datasets, but their advantage must be demonstrated rather than assumed. A fair assessment requires that they be tested against evaluated deep-learning baselines under the same folds, metrics, and model-selection constraints.

Prior grinding roughness studies have often lacked at least one of the following elements: strict condition-level validation, systematic comparison of signal representations, auditable and deep models evaluated under identical folds, statistical testing at the condition level, seed sensitivity analysis, uncertainty calibration under grouped condition shift, condition-level failure analysis, or latency and model-footprint comparison. As a result, it remains unclear whether physics-motivated and more auditable models can provide competitive \(R_a\) prediction when the test condition is entirely unseen. The central question addressed in this work is therefore: can directly inspectable linear models and tree ensembles supported by post-hoc attribution provide competitive \(R_a\) prediction relative to evaluated neural baselines under strict leave-one-condition-out validation, while exposing uncertainty or failure modes under grouped condition shift?

This study investigates in-process \(R_a\) prediction from multimodal AE, vibration, and process-parameter data collected during grinding. The dataset contains 319 valid grinding passes grouped into 16 grinding conditions generated by a fractional 16-condition design; it is not a full \(4 \times 4 \times 4\) factorial experiment. Generalization is assessed by strict leave-one-condition-out cross-validation with 16 folds, so each test fold corresponds to a previously unseen grinding condition. The benchmark inventory contains 25 exactly reconstructable model variants and two archival variants with incomplete resolved-configuration provenance across 38 input configurations, covering process parameters, time-domain statistics, handcrafted and physics-informed features, dB spectrograms, dB-z descriptors, and two-dimensional wavelet-scattering representations for AE and vibration. A directly inspectable ridge baseline, auditable tree ensembles with post-hoc attribution, and a separate shallow neural baseline are compared with evaluated deep-learning architectures, including ResNet-style spectrogram CNNs, multimodal fusion CNNs, bilinear fusion, trajectory CNN, graph-fusion, TabTransformer, and attention-based regressors. Explainability is examined using ridge coefficients, TreeSHAP, and Grad-CAM, but the resulting attributions are interpreted as physically plausible associations rather than causal proof. Uncertainty is evaluated under full leave-one-condition-out validation using MC-dropout and RF OOB/tree-variance intervals, with emphasis on under-coverage and condition-level outliers rather than on solved calibration. Finally, feature-precomputed model-inference timing and model-footprint analyses are reported separately from full end-to-end deployment latency, which would also require the acquisition window, STFT/dB-z/log-mel computation, feature scaling, communication overhead, and edge-hardware testing to be considered.

In the present benchmark, random forests trained on pass-level fused AE and vibration dB-z or log-mel descriptors form the leading observed point-estimate group. This result should be interpreted cautiously. The primary planned statistical comparison establishes conclusive superiority only over the shallow MLP baseline, while a current fixed-setting five-seed analysis gives the same qualitative ordering for selected baselines but does not reproduce the historical tuned checkpoints. In addition, the best MAE is close to the manufacturer-quoted repeatability of the profilometer, and raw trace-level repeatability or gauge R\&R data are unavailable. Consequently, sub-\(0.005\)--\(0.010~\upmu\)m ranking differences are not treated as practically meaningful. The intended contribution is therefore not a general rejection of deep learning, but a condition-level, representation-aware, and measurement-conscious comparison of physics-aware machine-learning strategies for interpretable grinding roughness prediction.
